%% Abstract placeholder

Unlike C, C++, or even the most recent versions of Java, Standard ML does not
expose direct control of memory layouts to the programmer: it is impossible to
observe in the source language whether an object of an aggregate type is passed
as a boxed reference or an unboxed value. Consequently, the compiler has a wide
degree of freedom to choose representations so as to maximize performance. As a
general rule, on modern CPU architectures, the overarching goal of such choices
is to eliminate indirections and improve locality.

This paper explores two novel flattening optimizations: \texttt{PreFlatten},
which seeks to unbox aggregates passed across function calls,
and \texttt{ShallowFlatten}, which seeks to unbox aggregates stored in
containers. We implement thes optimizations in both the MLton compiler and its
parallel variant, MaPLe, and measure their impact across two
benchmark suites: MLton's (sequential) benchmarks, and
MaPLe's \texttt{parallel-ml-bench}.

Under ``fair'' conditions, we find that \texttt{PreFlatten} is broadly
neutral-to-negative (a geomean 2.7\% win to a 4.2\% loss across different
benchmark configurations) and \texttt{ShallowFlatten} is a moderate win (no
impacted sequential bencnhmarks; geomean wins of 13\% to 3\% on parallel
benchmarks). After controlling for known compiler and runtime artifacts
(for \texttt{PreFlatten}: placement of ``chunk'' boundaries in the generated
code; for \texttt{ShallowFlatten}: timing of GC pauses), the regressions
on \texttt{PreFlatten} are resolved (geomean 2.9\% win to 1.0\% loss),
and \texttt{ShallowFlatten} performance strongly improves (geomean wins of up to
50\%).

In the course of the evaluation, we also note that the default C toolchain
settings for both MLton and MaPLe lead to severe instability in performance
measurements, yielding per-benchmark noise of up to \hl{TODO: measure} against
no-op changes.

%% \hl{TODO: FILL IN}

%% We evaluate their impact across a suite of sequential and parallel
%% benchmarks, demonstrating performance wins of up to 50\%.
